% 02-architecture.tex --- Four-Pillar Decomposition (expanded to 10+ pages)
\section{四支柱分解}
\label{sec:arch}

该框架分为四个物理隔离的共享库和一个 Python 扩展，所有这些都由一次 CMake 调用生成。隔离不是社论；它由目录、目标、链接边缘和标头可见性强制执行。新手读\texttt{p10}来理解张量存储，无需学习任何微分概念；内核作者在接触算术内核后仅重建 \texttt{p10} ；编译器作者在 \texttt{stax} 内进行实验，无需重新链接存储层。成本是明确的边界和缺乏整个程序内联，这是该项目所接受的。

\subsection{目录 \texorpdfstring{->}{->} 工件映射}
\label{sec:arch:mapping}

每个支柱都是一个 CMake 目标，其 \texttt{OUTPUT\_DIRECTORY} 固定到 \texttt{tensorplay} 包树内的某个位置，因此树内构建无需安装即可导入，而轮子无需 \texttt{LD\_LIBRARY\_PATH} 即可导入。表~\ref{tab:pillars}为规范图。

\begin{table}[H]
\centering
\caption{目录到工件的映射。每行都是一个 CMake 目标； \texttt{OUTPUT\_DIRECTORY} 通过 \texttt{TP\_PACKAGE\_DIR} 是绝对的。}
\label{tab:pillars}
\small
\begin{tabular}{p{1.8cm}p{1.6cm}p{2.8cm}p{3.4cm}p{2.6cm}}
\toprule
目录 & CMake 目标 & 神器 (\texttt{OUTPUT\_DIRECTORY}) & 链接部门 & 责任 \\
\midrule
  & \texttt{p10} & \texttt{libp10.so} & \texttt{Threads}、\texttt{COMPUTE\_LIBS}（BLAS/oneDNN）、\texttt{CUDA::*} & 张量、存储、调度表、300 多个内核、分配器 \\
  & \texttt{tpx} & \texttt{libtpx.so} & \texttt{p10} & 显式 DAG、引擎、保存变量、衍生节点 \\
  & \texttt{stax} & \texttt{libstax.so} & \texttt{p10}, \texttt{tpx} & 最小 \texttt{Graph} IR、\texttt{Pass} 注册表、解释器 \\
  & \texttt{tp\_python} & \texttt{libtp\_python.so} & \texttt{p10}, \texttt{pybind11} & \texttt{CPythonBridge.cpp} 单拷贝脚轮基板 \\
  & \texttt{\_C} & \texttt{\_\_init\_\_.so} & \texttt{p10}, \texttt{tpx}, \texttt{stax}, \texttt{tp\_python} & 生成的操作、设备绑定、分析器、分布式 \\
\bottomrule
\end{tabular}
\end{table}

具体来说，\texttt{TP\_PACKAGE\_DIR} 默认为源代码树的 \texttt{tensorplay/} 目录，\texttt{TP\_LIB\_DIR} 为 \texttt{\$\{TP\_PACKAGE\_DIR\}/lib}。四个 \texttt{set\_target\_properties} 块强制将 \texttt{OUTPUT\_NAME "\_\_init\_\_"} 扩展至 \texttt{tensorplay/\_C}，并强制将 \texttt{p10}、\texttt{tpx}、\texttt{stax}、\texttt{tp\_python} 扩展至 \texttt{lib}。

Python 表面包含数百个模块，但只导入一个已编译的扩展。该单板重新导出通过枚举 \texttt{dir(\_C)} 和过滤 \texttt{\_\_module\_\_} 重写 (\S\ref{sec:veneer}) 发现的 550 多个本机符号。没有 Python 文件添加新的 CMake 目标；每个Python模块都是一个薄包装器，其繁重的工作重新进入\texttt{\_C}。

\begin{table}[H]
\centering
\caption{工件大小和链接可见性。在使用 \texttt{USE\_CUDA=ON}、\texttt{USE\_MKL=ON} 的发布版本上测量。}
\label{tab:artifactsizes}
\small
\begin{tabular}{lll}
\toprule
目标 & 链接可见性 & Why \\
\midrule
\texttt{p10} & \texttt{PUBLIC Threads::Threads}, \texttt{PUBLIC COMPUTE\_LIBS} & 线程池和 BLAS 符号必须对 \texttt{tpx}、\texttt{stax}、\texttt{\_C} 使用者可见 \\
\texttt{tpx} & \texttt{PRIVATE p10} 在第 ~26 行，\texttt{PUBLIC p10} 在第 ~40 行 & 手写节点需要私有p10头；下游 \texttt{stax} 传递继承 p10 接口 \\
\texttt{stax} & \texttt{PRIVATE tpx p10} 在第 ~8 行，\texttt{PUBLIC p10} 在第 ~17 行 & 解释器调用 \texttt{tpx::ops} 但仅公开 p10 张量类型 \\
\texttt{tp\_python} & \texttt{PUBLIC p10 pybind11::pybind11} & 单脚轮基板必须在 \texttt{\_C} 和树外操作模块之间共享 \\
\texttt{\_C} & \texttt{PRIVATE p10 tpx stax tp\_python tensorpipe} & 该扩展是将所有三个支柱连接在一起的唯一消费者 \\
\bottomrule
\end{tabular}
\end{table}

\subsection{构建顺序和空源记录}
\label{sec:arch:buildorder}

仅用四行声明顺序：

\begin{lstlisting}[caption={pillar order is load-bearing.}, label=lst:buildorder]
# Add p10
add_subdirectory(p10)          # :938
# Add tpx
add_subdirectory(tpx)          # :941
# Add stax (Optional)
add_subdirectory(stax)         # :944
\end{lstlisting}

私人链接 \texttt{p10} ；公开重新导出它，以便依赖 \texttt{tpx} 的消费者自动获得 p10 接口。私下链接 \texttt{p10} 和 \texttt{tpx} 并在第 17 行公开重新导出 \texttt{p10}。 \texttt{\_C} 与 \texttt{p10 tpx stax tp\_python} 的链接：这是三个级别相交的唯一地方。对这三个 \texttt{add\_subdirectory} 调用中的任何一个重新排序或添加后沿都会立即失败，并出现 CMake 配置错误或链接时未定义的引用，而不是静默运行时错误顺序。

生成步骤先于所有这些。在 \texttt{custom\_command} 处，通过 \texttt{main.py} 从 \texttt{native\_functions.yaml} 和 \texttt{derivatives.yaml} 生成十个工件。每个支柱目标都声明 \texttt{add\_dependencies(generate\_code)}（例如\,）。生成的 header 通过 \texttt{TensorGenerated.h"} 混入；生成的源作为 \texttt{\$\{GEN\_CPP\}} 附加到 \texttt{P10\_SOURCES} 处。因此，YAML 编辑仅重建更改后的生成文件及其依赖文件。

保留单个注释行：

\begin{lstlisting}[caption={The vacated-source comment in the p10 build file: metadata that once violated the axiom was physically moved out, leaving this trail.}, label=lst:audittrail]
    # src/AutogradMetaInterface.cpp (Moved to TPX)
\end{lstlisting}

这是故意的。早期版本在 \texttt{p10} 中实现了具体元数据；当前的设计使 \texttt{p10} 头与 \texttt{tpx} 无关，因此 USB 插槽就足够了。将注释留在源列表中会记录决策边界：grep 注释证明该文件曾经存在于 \texttt{p10} 中，并且被故意腾空以保持非循环性。删除它会抹去为什么 \texttt{p10} 不能包含 \texttt{tpx} 的历史记录；在 \texttt{p10} 内复活它会重新引入后沿并破坏必须返回空的检查的不变量。

出现了相关学科。在 Windows 上，对于 Linux 轮子，\texttt{BUILD\_SHARED\_LIBS} 通常是 \texttt{ON}，块强制 \texttt{BUILD\_SHARED\_LIBS OFF}：

\begin{lstlisting}[caption={Windows static embedding.}, label=lst:winstatic]
if(WIN32 AND BUILD_SHARED_LIBS)
    set(BUILD_SHARED_LIBS OFF)
endif()
\end{lstlisting}

如果没有这个开关，每个 DLL 都需要在每个跨越边界的 p10 符号上使用 \texttt{\_\_declspec(dllexport)} ，而代码库并不携带该符号。 Windows 上的静态链接将三个支柱嵌入到单个模块 \texttt{\_C.pyd} 中，并保持符号可见性与 Linux 共享库布局相同，而无需注释每个类。

\subsection{产品布局：RPATH 和两个预加载路径}
\label{sec:arch:rpath}

构建后，包树就是产品：

\begin{lstlisting}[caption={Post-build layout under the package directory.}, label=lst:layout]
tensorplay/
  lib/
    libp10.so
    libtpx.so
    libstax.so
    libtp_python.so
    cudart64_*.dll / cublas64_*.dll / cudnn*.dll   # Windows wheels only
  _C/
    __init__.so          # OUTPUT_NAME "__init__" :1176
    __init__.pyi
  __init__.py            # 1,197 lines, three dirty jobs
  _tensor.py             # Tensor = _C.TensorBase :12
  version.py             # generated at :330
\end{lstlisting}

相对 RPATH 使布局可重定位。从：

\begin{lstlisting}[caption={RPATH is the loader contract.}, label=lst:rpath]
if(NOT WIN32)
    set_target_properties(p10 tpx stax PROPERTIES
        INSTALL_RPATH "$ORIGIN"
        BUILD_RPATH "$ORIGIN"
    )
    set_target_properties(tp_python PROPERTIES
        INSTALL_RPATH "$ORIGIN"
        BUILD_RPATH "$ORIGIN"
    )
    set_target_properties(_C PROPERTIES
        INSTALL_RPATH "$ORIGIN/../lib"
        BUILD_RPATH "$ORIGIN/../lib"
    )
endif()
\end{lstlisting}

\texttt{p10}、\texttt{tpx}、\texttt{stax} 和 \texttt{tp\_python} 一起驻留在 \texttt{lib} 中，因此 \texttt{\$ORIGIN} 就足够了：每个 \ac{dso} 都可以在没有绝对路径的情况下找到其兄弟姐妹，并且当树移动时不需要 \ac{rpath} 编辑。 \texttt{\_C} 位于 \texttt{\_C} 更深的一个目录中，因此它使用 \texttt{lib}。验证是一个不涉及 Python 的命令：

\begin{lstlisting}[caption={Linux loader verification without importing Python.}, label=lst:readelf]
readelf -d tensorplay/lib/libp10.so | grep -i rpath      # $ORIGIN
readelf -d tensorplay/lib/libtpx.so | grep -i rpath      # $ORIGIN
readelf -d tensorplay/_C/__init__.so | grep -i rpath     # $ORIGIN/../lib
ldd tensorplay/_C/__init__.so | grep "not found"          # must be empty
\end{lstlisting}

当wheel安装到\texttt{LD\_LIBRARY\_PATH}为空的虚拟环境中时，没有此块的失败表现为\texttt{ImportError: libp10.so: cannot open shared object file}，这正是wheel的目标场景。的注释将其记录为 RPATH 存在的唯一原因； Windows 完全忽略 RPATH，因此 Windows 分支由 \texttt{if(NOT WIN32)} 保护。

Windows 使用显式预加载。定义由 \texttt{sys.platform == 'win32'} 保护的 \texttt{\_load\_dll\_libraries}：

\begin{lstlisting}[caption={ordered DLL preload on Windows.}, label=lst:winpreload]
kernel32 = ctypes.WinDLL("kernel32.dll", use_last_error=True)
# MKL -> CUDA/cuDNN -> p10/tpx/stax   :80
mkl_dlls = glob.glob(os.path.join(package_lib_path, "mkl_*.dll"))
cuda_dlls = glob.glob(os.path.join(package_lib_path, "cudart*.dll")) + \
            glob.glob(os.path.join(package_lib_path, "cublas*.dll")) + \
            glob.glob(os.path.join(package_lib_path, "cudnn*.dll"))
core_dlls = [os.path.join(package_lib_path, "p10.dll"),
             os.path.join(package_lib_path, "tpx.dll"),
             os.path.join(package_lib_path, "stax.dll")]
for dll in sorted_mkl + cuda_dlls + core_dlls:
    if with_load_library_flags:
        res = kernel32.LoadLibraryExW(dll, None, 0x00001100)  # :122
    else:
        res = kernel32.LoadLibraryW(dll)
\end{lstlisting}

\texttt{0x00001100} 是 \texttt{LOAD\_LIBRARY\_SEARCH\_DEFAULT\_DIRS | LOAD\_LIBRARY\_SEARCH\_DLL\_LOAD\_DIR}，它仍然搜索之前通过 \texttt{os.add\_dll\_directory} 添加的包 lib 目录的最低许可标志。顺序是承载的：\texttt{mkl\_core} 和 \texttt{mkl\_sequential} 在 MKL 组内首先在第 85--94 行排序，因为 \texttt{mkl\_def} 和 \texttt{mkl\_avx2} 依赖于它们； CUDA 运行时在 \texttt{p10} 之前加载，因此 p10 的 CUDA 内核所需的 \texttt{CUDA::cudart} 符号已经解析。帮助程序 \texttt{\_add\_dll\_directory} 尝试 \texttt{lib} 和旧根，加上解释器的 \texttt{bin}、基本环境和用户站点包 bin，每个都由 \texttt{os.path.exists} 保护，因此丢失的 Conda 安装不会出错。

在 Linux 上，对应部分不会在仅 CPU 环境中急切地加载任何内容。包装扩展导入：

\begin{lstlisting}[caption={Linux CUDA wheel probing.}, label=lst:linuxpreload]
try:
    from . import _C as _C                                          # :199
except (ImportError, OSError) as _load_error:                         # :200
    if sys.platform.startswith('linux') and '_preload_cuda_deps' in globals():
        _preload_cuda_deps(_load_error)                              # :202
        from . import _C as _C                                       # :203
    else:
        raise

def _preload_cuda_deps(err=None):                                     # :174
    cuda_error = any(token in str(err) for token in
                     ('libcud', 'libcublas', 'libcurand', 'libnvrtc'))  # :177
    if not cuda_error:
        raise err
    for lib_folder, lib_name, required in (                           # :185
        ('cublas', 'libcublasLt.so.*[0-9]', True),
        ('cublas', 'libcublas.so.*[0-9]', True),
        ('cudnn',  'libcudnn.so.*[0-9]', True),
        ('cuda_nvrtc', 'libnvrtc.so.*[0-9]', False),
        ('cuda_runtime','libcudart.so.*[0-9]', True),
        ('curand', 'libcurand.so.*[0-9]', True),
    ):
        _preload_cuda_lib(lib_folder, lib_name, required=required)   # :193
\end{lstlisting}

\texttt{\_get\_cuda\_dep\_paths} 位于第 144 行，每个 \texttt{sys.path} 条目内的探头 \texttt{libcublasLt.so.*} 和 \texttt{libcudart.so.*}，这正是点轮 \texttt{nvidia-cublas-cu*} 和 \texttt{nvidia-cuda-runtime-cu*} 安装的布局。如果探测失败，它将回退到 \texttt{libcublasLt.so.*}，即捆绑工具包时使用的路径。然后执行 \texttt{ctypes.CDLL(candidates[0])}，它在解析扩展的 \texttt{DT\_NEEDED} 之前加载 SONAME，解决用户将 CUDA 轮安装到环境的 \texttt{site-packages} 而不设置 \texttt{LD\_LIBRARY\_PATH} 的常见情况。仅拦截提及 \texttt{libcud}、\texttt{libcublas}、\texttt{libcurand} 或 \texttt{libnvrtc} 的错误；所有其他 \texttt{ImportError} 原因均按原样重新提出。

\begin{table}[H]
\centering
\caption{加载程序行为：Linux RPATH vs.\ Windows 预加载 vs.\ Linux CDLL 探针。}
\label{tab:loader}
\small
\begin{tabular}{llll}
\toprule
平台 & 机制 & 文件 & 引文 \\
\midrule
Linux（全部） & \texttt{INSTALL\_RPATH "\$ORIGIN"} & \texttt{libp10.so} 兄弟姐妹 & \\
Linux（扩展） & \texttt{lib"} & \texttt{\_\_init\_\_.so} & \\
视窗 & \texttt{LoadLibraryExW} 与 \texttt{0x00001100} & \texttt{stax.dll} & \\
Linux+CUDA轮子 & \texttt{lib} 的 \texttt{ctypes.CDLL} 探针 & \texttt{libcublasLt.so.*} & \texttt{\_\_init\_\_.py:161,186} \\
捆绑 CUDA & \texttt{CTOOLKIT\_BIN\_DIR} 全局 + \texttt{install(FILES)} & \texttt{cublasLt64\_*.dll} & \\
\bottomrule
\end{tabular}
\end{table}

\subsection{虚拟接口解耦：USB 插槽}
\label{sec:arch:usb}

不能包含 \texttt{tpx} 但必须携带需要梯度的张量的微分状态。天真的替代方案都被拒绝了：将具体字段（\texttt{grad\_fn\_}，\texttt{grad\_}）放入\texttt{p10}中，将存储与微分结合起来；让 \texttt{p10} 包含 \texttt{tpx} 标头会创建一个循环依赖关系，\texttt{cmake --build} 无法进行拓扑排序，并且会破坏增量构建。相反，该设计在 \texttt{p10} 中声明一个纯抽象槽，并插入从 \texttt{tpx} 到 \texttt{shared\_ptr<AutogradMetaBase>} 的具体对象以及 \texttt{static\_cast} 握手。

该插槽有四个操作，每个操作都是纯虚拟的：

\begin{lstlisting}[caption={the pure interface.}, label=lst:autogradbase2]
class P10_API AutogradMetaBase {
public:
  virtual ~AutogradMetaBase() = default;
  virtual bool requires_grad() const = 0;
  virtual void set_requires_grad(bool) = 0;
  virtual Tensor grad() const = 0;
  virtual void set_grad(const Tensor&) = 0;
  virtual bool retains_grad() const = 0;
  virtual void set_retains_grad(bool) = 0;
};
\end{lstlisting}

标头仅包含 \texttt{Macros.h}；它转发 \texttt{class Tensor} 而不包括 \texttt{Tensor.h}，因此编译时占用空间是单个 vtable 加上 6 个虚拟槽（x86-64 上为 48 字节）。没有内核、没有引擎、没有图形概念跨越这里的边界。

持有插槽：

\begin{lstlisting}[caption={never-copied owner.}, label=lst:implslot]
  // Autograd extension point.
  // Never copied by TensorImpl copy operations: copies start without history.
  std::shared_ptr<AutogradMetaBase> autograd_meta_;
\end{lstlisting}

at 的复制构造函数故意省略了该成员：

\begin{lstlisting}[caption={history is not duplicated.}, label=lst:copyctor]
TensorImpl::TensorImpl(const TensorImpl& other)
    : storage_offset_(other.storage_offset_),
      sizes_and_strides_(other.sizes_and_strides_),
      dtype_(other.dtype_),
      device_(other.device_),
      version_counter_(other.version_counter_),
      inference_tensor_(other.inference_tensor_),
      is_contiguous_(other.is_contiguous_),
      memory_format_(other.memory_format_),
      storage_(other.storage_),
      shared_state_(other.shared_state_),
      sparse_state_(other.sparse_state_),
      transform_value_(other.transform_value_),
      transform_batch_dim_(other.transform_batch_dim_),
      transform_level_(other.transform_level_),
      quantizer_(other.quantizer_) {}
      // autograd_meta_ is intentionally absent: default-initialized to nullptr
\end{lstlisting}

因此 \texttt{Tensor b = a;} 共享存储、大小、步幅、设备、数据类型、稀疏状态、量化器和版本计数器别名，但不共享图。不变量是：形状突变和存储别名传播；历史没有。移动操作是默认的，\texttt{copy\_metadata\_from} at 保留了 \texttt{autograd\_meta\_} 的省略。
这与“从不复制”合约相匹配，并启用 \texttt{auto y = x + 1; y.backward}，其中 \texttt{x} 保持叶子状态，而 \texttt{y} 成为其子节点，而无需别名元数据。

混凝土塞位于：

\begin{lstlisting}[caption={concrete fields, tpx only.}, label=lst:autogradconcrete]
class TENSORPLAY_API AutogradMeta : public AutogradMetaBase {
private:
  bool requires_grad_ = false;
  bool retains_grad_ = false;
  tensorplay::Tensor grad_;
  std::shared_ptr<Node> grad_fn_;
  std::weak_ptr<Node> grad_accumulator_; // avoids leaf cycle
  uint32_t output_nr_ = 0;
  bool has_view_info_ = false;
  tensorplay::Tensor view_base_;
  std::vector<int64_t> view_sizes_;
  std::vector<int64_t> view_strides_;
  int64_t view_storage_offset_ = 0;
  std::function<tensorplay::Tensor(const tensorplay::Tensor&)> view_fn_;
  mutable std::mutex view_mutex_;
};
\end{lstlisting}

\texttt{grad\_accumulator\_} 是 \texttt{weak\_ptr<Node>} 而不是 \texttt{shared\_ptr<Node>}，因为 \texttt{AccumulateGrad} 强烈拥有其叶张量；强大的后边缘将形成一个无法收集的循环，该循环会泄漏参与图中的每个叶子。所有其他成员均拥有强大股权； \texttt{view\_fn\_} 加上 \texttt{view\_mutex\_} 支持陈旧视图重播。

握手是两个助手，每个助手都有一个 \texttt{static\_cast}，热路径上没有 \texttt{dynamic\_cast}：

\begin{lstlisting}[caption={the cast boundary.}, label=lst:handshake]
AutogradMeta* get_autograd_meta(const Tensor& t) {
    if (auto* impl = t.unsafeGetTensorImpl().get()) {
        return static_cast<AutogradMeta*>(impl->autograd_meta()); // :17
    }
    return nullptr;
}
AutogradMeta* get_or_create_autograd_meta(const Tensor& t) {
    auto impl = t.unsafeGetTensorImpl();
    if (!impl) return nullptr;
    if (auto* meta = impl->autograd_meta()) {
        return static_cast<AutogradMeta*>(meta);                  // :26
    }
    auto meta = std::make_shared<AutogradMeta>();
    auto* raw = meta.get();
    impl->set_autograd_meta(std::move(meta));                      // :30
    return raw;
}
\end{lstlisting}

第 17 行的调用是安全的，因为分配给 \texttt{autograd\_meta\_} 的唯一具体子类是 \texttt{AutogradMeta}，它在第 28 行构造并在第 30 行安装。 \texttt{p10} 中的标头不会看到 \texttt{AutogradMeta.h}，因此构建不会列出包含路径。静态类型系统加上共享库边界共同强制执行不变性。
第二个证明是当转换目标更改时构建失败：如果 \texttt{tpx} 重命名 \texttt{AutogradMeta} 而不更新转换，则扩展 \texttt{\_C} 将无法与 \texttt{undefined reference to tensorplay::tpx::AutogradMeta} 链接，而不是在运行时损坏张量。

此设计的每次调用成本是 \texttt{requires\_grad} / \texttt{grad} / \texttt{set\_grad} 的一次间接虚拟调用加上 \texttt{autograd\_meta\_} 的一次指针间接调用。好处是，当内核更改不再触发 \texttt{tpx} 重建时，\texttt{p10} 会重建，而当引擎调度更改不再触发 \texttt{p10} 重建时，\texttt{tpx} 会重建。

\begin{tpdef}[Definition (USB Slot Pattern)]
\emph{USB slot} 是在较低级别库 $L$ 中声明的虚拟接口 $I$，加上嵌入在 $L$ 状态中的 \texttt{shared\_ptr<I>} 类型的不透明拥有指针，其中具体实现 $C{:}I$ 仅存在于较高级别库 $H$ 中。该模式是健全的 iff (i) $L$ 从不包含任何 $H$ 标头； (ii) $L$ 的复制构造函数中省略了指针（复制开始时没有历史记录）； (iii) $H$ 边界向下是 \texttt{static\_cast}，这是合理的，因为 $C$ 是唯一安装到插槽中的类。 TensorPlay 的实例：$I{=}$\texttt{AutogradMetaBase}、$L{=}$\texttt{p10}、$C{=}$\texttt{AutogradMeta}、$H{=}$\texttt{tpx}。
\end{tpdef}

\subsection{图：跨共享库边界的 USB 插槽模型}
\label{sec:arch:figure}

图~\ref{fig:usb}呈现了与硬件类比相同的机制：插槽在下部组件中标准化，插头由上部组件提供，电缆是虚拟调度。虚线是\texttt{ldd}报告的\texttt{.so}边界；它左边的符号在\texttt{libp10.so}中，它右边的符号在\texttt{libtpx.so}中。该图由 \texttt{gen\_01\_slot.py} 在本地复制，但事实来源是标题中引用的两个标题。

\begin{figure}[H]
\centering
\begin{tikzpicture}[node distance=0.55cm, every node/.style={font=\footnotesize}]
\node[libbox, fill=blue!8, draw=tpblue, minimum width=5.4cm, minimum height=2.6cm, align=left] (p10b) at (0,0) {\textbf{p10 $\cdot$ TensorImpl}\\ \texttt{shared\_ptr<AutogradMetaBase> autograd\_meta\_}\\ \texttt{ never copied (TensorImpl.cpp:77)}\\ \textit{USB slot (pure virtual)} \\ \texttt{ set\_grad}};
\node[libbox, fill=orange!8, draw=tporange, minimum width=5.4cm, minimum height=2.6cm, align=left, right=1.35cm of p10b] (tpxb) {\textbf{tpx $\cdot$ AutogradMeta}\\ \texttt{grad\_fn\_: shared\_ptr<Node>}\\ \texttt{grad\_accumulator\_: weak\_ptr<Node>}\\ \texttt{ view\_mutex\_}\\ \textit{USB plug (concrete production object)}};
\draw[arrow, color=tporange, thick] ([yshift=0.55cm]p10b.east) -- node[above, font=\scriptsize\ttfamily, fill=white, inner sep=1pt] {static\_cast<AutogradMeta*>} ([yshift=0.55cm]tpxb.west);
\draw[arrow, color=tpblue, thick, dashed] ([yshift=-0.55cm]tpxb.west) -- node[below, font=\scriptsize\ttfamily, fill=white, inner sep=1pt] {set\_autograd\_meta(move(meta))} ([yshift=-0.55cm]p10b.east);
\draw[dashed, color=gray, line width=0.6pt] (2.72,-1.7) -- (2.72,1.7);
\node[font=\scriptsize\itshape, color=gray, fill=white, inner sep=1pt] at (2.72,1.95) {shared-lib boundary: \texttt{\_\_init\_\_.so}};
\node[font=\scriptsize\itshape, color=gray] at (2.72,-1.95) {no \texttt{..."} in \texttt{p10} $\;\Rightarrow\;$  $=0$};
% small vtable hint
\node[font=\tiny\ttfamily, color=tpblue, below=0.05cm of p10b, anchor=north] {vtable @ AutogradMetaBase.h:13};
\node[font=\tiny\ttfamily, color=tporange, below=0.05cm of tpxb, anchor=north] {overrides @ AutogradMeta.h:52};
\end{tikzpicture}
\caption{USB 插槽型号。 \texttt{p10} 声明抽象槽； \texttt{tpx} 制造插头；握手时间为 \texttt{static\_cast}，在 \texttt{Autograd.cpp}。 at 的复制构造函数省略了槽。}
\label{fig:usb}
\end{figure}

\subsection{设计公理：层次单调性}
\label{sec:arch:axiom}

令 $V$ 为构建目标集，并令 $D\subseteq V\times V$ 为关系“$a$ 链接 $b$”（即\ $a$ 与 $b$ 有 \texttt{target\_link\_libraries} 边缘）。定义 \texttt{level}: $V\rightarrow \mathbb{N}$ 通过

\begin{table}[H]
\centering
\caption{级别为总级别并且等于 \texttt{add\_subdirectory} 顺序。}
\label{tab:levels}
\small
\begin{tabular}{clcl}
\toprule
等级 & 目标 & \texttt{add\_subdirectory} 线 & 允许从更高层导入标头 \\
\midrule
0 & \texttt{p10} &  & 不存在更低的目标 \\
1 & \texttt{tpx} &  & 可能包括 \texttt{*.h} \\
2 & \texttt{stax} &  & 可能包括 \texttt{include}、\texttt{include} \\
3 & \texttt{\_C} / \texttt{tp\_python} & , \texttt{1106} & 可能包括 \texttt{p10}、\texttt{tpx}、\texttt{stax} \\
\bottomrule
\end{tabular}
\end{table}

公理（水平单调）：$\forall (a,b)\in D,\; \mathrm{level}(a) > \mathrm{level}(b)$。

同样，每个边都严格按照 \texttt{p10=0 < tpx=1 < stax=2 < \_C=3} 的顺序向下指向。表~\ref{tab:levels}表明\texttt{D}仅包含\texttt{(tpx,p10)}、\texttt{(stax,p10)}、\texttt{(stax,tpx)}、\texttt{(tp\_python,p10)}、\texttt{(\_C,p10)}、\texttt{(\_C,tpx)}、\texttt{(\_C,stax)}、\texttt{(\_C,tp\_python)}，每个都满足严格的不等式。不存在与 \texttt{level}(a) $\le$ \texttt{level}(b) 的配对；特别是 \texttt{(p10,tpx)} 和 \texttt{(p10,stax)} 不存在。

强制执行是由链接器执行的，而不是按照惯例。在任何 \texttt{Tensor*.h} 中添加 \texttt{AutogradMeta.h"} 会引入对 \texttt{tensorplay::tpx::AutogradMeta::requires\_grad} 的引用，该 \texttt{tensorplay::tpx::AutogradMeta::requires\_grad} 的定义位于 \texttt{libtpx.so} 中。因为 \texttt{p10} 不链接 \texttt{tpx}（根据公理，也不能链接），所以 \texttt{libp10.so} 的链接会成功，但任何单独链接 \texttt{p10} 的使用者的链接都会因 \texttt{undefined reference to tensorplay::tpx::AutogradMeta} 而失败。在 Linux 上，这是在 \texttt{cmake --build} 期间诊断的；在 Windows 上，诊断为 \texttt{\_C.pyd} 的 \texttt{link.exe} 步骤。传统整体框架的替代故障模式是仅在程序启动时诊断的静态初始化顺序缺陷。

该公理还决定了重建粒度。触摸 \texttt{ActivationKernels.cpp} 只会弄脏 \texttt{p10} 和 \texttt{\_C} 的最后一个链接； \texttt{tpx} 和 \texttt{stax} 不会重建，因为它们不依赖于 \texttt{p10} 源，仅依赖于其安装的标头。触摸 \texttt{Engine.cpp} 会重建 \texttt{tpx} 并重新链接 \texttt{stax} 和 \texttt{\_C}，但不会重新链接 \texttt{p10}。 
一个推论涉及\texttt{tp\_python}。它是一个 \texttt{SHARED} 库 \texttt{CPythonBridge.cpp}，仅链接 \texttt{p10} 和 \texttt{pybind11}。它的存在正是出于一个原因，记录了 pybind11 Caster 基底和包装身份缓存必须从单个模块边界注册。如果 \texttt{CPythonBridge.cpp} 在 \texttt{p10} 和 \texttt{\_C} 中分别编译，则每个副本都将携带自己的注册表，并且返回到 Python 的第一个 \texttt{Tensor} 将出现段错误。第~3级的单个共享副本保留\texttt{level(tp\_python)=3 > level(p10)=0}，并且仍然允许树外操作模块通过\texttt{add\_tensorplay\_op}链接相同的底层，而无需重新编译核心。

\subsection{Python Veneer：1,197 行中的三项肮脏工作}
\label{sec:veneer}

\texttt{\_\_init\_\_.py} 有 1,197 行，与 \texttt{\_tensor.py} 一起构成了编译后的扩展无法表达的整个 Python 级别的设计。所有其他 Python 模块（\texttt{nn}、\texttt{optim}、\texttt{autograd}、\texttt{functional}、\texttt{distributed}）最终都会重新进入 \texttt{\_C}。三个问题不可能在 C++ 中干净地实现，因此集中在这里。

\subsubsection{作业 1：DLL/共享对象预加载（第 34--205 行）}
如 \S\ref{sec:arch:rpath} 中所述，此作业处理 Linux 加载程序（尊重 \texttt{RPATH}）和 Windows 加载程序（忽略它）之间的差异。在 Windows 上，该工作是强制性的；在 Linux 上，这是对 CUDA 轮子的机会性修复。代码在第 34 行和第 143 行分开，没有重叠，因此读者可以检查一个平台而无需了解另一个平台。该工作是唯一导入 \texttt{ctypes} 和 \texttt{glob} 的地方；也没有内核或引擎文件引用。

\subsubsection{作业 2：符号重命名 (\texttt{\_C.add} $\rightarrow$ \texttt{tensorplay.add})}

该扩展将每个运算符的 C++ 名称导出为 \texttt{\_C.add}、\texttt{\_C.relu}、\texttt{\_C.conv2d}。胶合板将它们重新导出为平面 \texttt{tensorplay.*} 名称并重写 \texttt{\_\_module\_\_}，以便回溯和 \texttt{help(tensorplay.add)} 报告公共位置。循环是：

\begin{lstlisting}[caption={renaming with builtin fallback.}, label=lst:renaming]
__attr_name, __obj = "", None
for __attr_name in dir(_C):
    if __attr_name[0] != "_" and not __attr_name.endswith("Base"):
        __all__.append(__attr_name)
        __obj = getattr(_C, __attr_name)
        if callable(__obj) or inspect.isclass(__obj):
            if os.getenv("TENSORPLAY_BUILDING_STUBS") == "1":
                continue
            if __obj.__module__ != __name__:
                try:
                    __obj.__module__ = __name__          # pybind11 class: succeeds
                except AttributeError:                    # builtin_function_or_method: fails
                    if callable(__obj):
                        def make_wrapper(f):
                            @functools.wraps(f)
                            def wrapper(*args, **kwargs):
                                return f(*args, **kwargs)
                            wrapper.__module__ = __name__
                            return wrapper
                        wrapper = make_wrapper(__obj)
                        if __attr_name in globals():
                            globals()[__attr_name] = wrapper
            if __attr_name not in globals():
                globals()[__attr_name] = __obj
\end{lstlisting}

\texttt{except} 存在是因为 pybind11 函数是 \texttt{builtin\_function\_or\_method} 描述符，其 \texttt{\_\_module\_\_} 属性在 CPython 中是只读的；后备包装器通过 \texttt{functools.wraps} 保留 \texttt{inspect.signature}，同时将 \texttt{\_\_module\_\_} 固定为 \texttt{"tensorplay"}。如果没有这项工作，用户将调用 \texttt{tensorplay.\_C.add} 或看到回溯提及 \texttt{tensorplay.\_C} 而不是 \texttt{tensorplay}。

张量类型本身在这里完成：定义 \texttt{Tensor = \_C.TensorBase} ，然后将方法（\texttt{t}、\texttt{type}、\texttt{cuda}、\texttt{cpu}、\texttt{long}、\texttt{float}，加上钩子、按位 dunders、\texttt{permute}、\texttt{expand}、\texttt{unique}、\texttt{topk}）附加到该别名上。 \texttt{Tensor.cpp} 处的 C++ 绑定在 \texttt{TensorBase} 上使用 \texttt{py::dynamic\_attr}，因此像 \texttt{class MyTensor(Tensor)} 这样的 Python 子类可以附加每个实例的属性，同时仍能被本机施法者识别。通过 \texttt{tp\_python} 中的共享施法者进行检查 \texttt{py::isinstance< Tensor >(obj)}。

\subsubsection{作业 3：通过 \texttt{\_\_getattr\_\_} 进行惰性子包（第 1118--1134 行）}

在单板顶层直接使用 \texttt{import tensorplay.linalg} 将强制在每个 \texttt{import tensorplay} 上导入 354 个模块，即使对于仅使用 \texttt{tensorplay.tensor([1])} 的脚本，导入时间也会增加数百毫秒。相反，胶合板注册了一个惰性集：

\begin{lstlisting}[caption={lazy subpackages.}, label=lst:lazy]
else:
    _lazy_modules = {
        "audio",
        "export",
        "func",
        "fft",
        "linalg",
        "quantization",
        "sparse",
        "special",
        "vision",
    }
    def __getattr__(name):
        if name in _lazy_modules:
            return importlib.import_module(f".{name}", __name__)
        raise AttributeError(f"module '{__name__}' has no attribute '{name}'")
\end{lstlisting}

\texttt{\_\_getattr\_\_} 是Python~3.7 中引入的模块级钩子。它仅在缺少属性时调用，因此 \texttt{tensorplay.Tensor} 或 \texttt{tensorplay.add} 等已知名称永远不会进入此路径。在 \texttt{TYPE\_CHECKING} 下（仅静态分析为 true，运行时为 false），胶合板改为在第 1108--1115 行急切地导入 \texttt{export} 和 \texttt{func}，以便 pylance 和其他 LSP 服务器解析属性而不执行惰性挂钩。随附的 \texttt{TYPE\_CHECKING} 导入由 \texttt{if TYPE\_CHECKING:} 保护，并且没有运行时成本。

\begin{table}[H]
\centering
\caption{三项肮脏的工作：地点、机制和被拒绝的单一替代方案。}
\label{tab:threedirty}
\small
\begin{tabular}{p{2.0cm}p{3.2cm}p{3.6cm}p{3.4cm}}
\toprule
Job & 线路 & 机制 & 被拒绝的替代方案 \\
\midrule
DLL预加载 &  & \texttt{LoadLibraryExW}, \texttt{ctypes.CDLL}, \texttt{add\_dll\_directory} & Windows \texttt{RPATH}（不存在） \\
符号重命名 &  & \texttt{\_\_module\_\_} 重写 + \texttt{functools.wraps} 包装器 & 在 C++ 中重命名（破坏 \texttt{dir(\_C)} 自省） \\
懒包 &  & \texttt{\_lazy\_modules} + \texttt{\_\_getattr\_\_} & Eager \texttt{import linalg} （增加导入延迟） \\
张量别名 & \texttt{Tensor = \_C.TensorBase} & \texttt{py::dynamic\_attr} 子类化 & 重复的 Python 类（破坏施法者身份） \\
\bottomrule
\end{tabular}
\end{table}

这三个工作加在一起是唯一具有跨领域影响的 Python 端复杂性。 \texttt{\_\_init\_\_.py} 中的其他内容都是声明性的： \texttt{DType} 别名位于第 230--259 行， \texttt{MemoryFormat} / \texttt{Layout} / \texttt{QScheme} 枚举位于第 264--305 行， \texttt{max}/\texttt{min}/\texttt{gradient} 包装器位于第 1038--1075 行，以及 \texttt{\_\_all\_\_} 扩展位于第 336--369 行枚举 550 多个具有排序合约的运算符。

\subsection{Stax：最小 IR，不是第二个前端}
\label{sec:arch:stax}

\texttt{stax} 是 CMake 的 76 行，用于从 \texttt{*.cpp} 生成 \texttt{libstax.so}。它的公共表面是两个标头：和。设计说明很明确：

\begin{lstlisting}[caption={not a second frontend.}, label=lst:staxpy]
"""Stax backend namespace.
Stax is not a second public compiler frontend.  Use ``tensorplay.compile``
and select ``backend="stax"`` once the backend has a production lowering.
The native ``tensorplay._C._stax`` module remains available for low-level IR
and optimizer development.
"""
from ._stax.stax import stax
\end{lstlisting}

将 \texttt{ValueNode} 定义为“生产者加槽”：

\begin{lstlisting}[caption={values are slots.}, label=lst:valuenode]
struct STAX_API ValueNode {
    size_t id;
    OpNode* producer; // Producer op, nullptr for graph inputs
    size_t producer_output_index; // Output index of producer
    std::string dtype = "float32";
    std::vector<int64_t> shape;
    std::string device = "cpu";
    std::vector<OpNode*> uses;
    ValueNode(size_t id_, OpNode* n, size_t off) : id(id_), producer(n), producer_output_index(off) {}
};
\end{lstlisting}

\texttt{OpNode} 通过第 62 行的 \texttt{variant<int64\_t, double, string, vector<int64\_t>, vector<double>>} 映射对输入、输出、所属图和属性进行分组。 \texttt{Graph} at 拥有 \texttt{nodes} 和 \texttt{values} 向量以及 \texttt{inputs} / \texttt{outputs} 边界，在第 95--97 行有辅助函数 \texttt{addInput}、\texttt{createNode} 和 \texttt{registerOutput}。 \texttt{IRBuilder} 帮助器创建命名节点并从第一个输入传播 dtype。

执行是严格顺序的。 \texttt{Graph::execute} 在第 192 行将具体的 \texttt{Tensor} 参数绑定到 \texttt{env[ input->id ]}，然后按创建顺序迭代 \texttt{nodes}。每个节点通过常规调度程序调度到 \texttt{tpx::ops} （例如\ \texttt{add}、\texttt{conv2d}、\texttt{matmul}、\texttt{max\_pool2d}、\texttt{batch\_norm}），保留设备和区分语义。中间生命周期由~198行的\texttt{remaining\_uses}和~199行的\texttt{keep\_alive}界定；死值通过第 215 行的 \texttt{env[id]=Tensor} 清除，允许分配器回收单个图形运行中的存储，而无需保留每个中间值，直到函数返回。

添加了唯一的优化：具有单个注册通道的通道基础结构，将 \texttt{mul} 和 \texttt{add} 融合到 \texttt{fused\_mul\_add} （\texttt{p10} 中的单个内核）。注册表是 \texttt{map<string, Creator>}，为了兼容而保留第 73 行的 \texttt{fuseGraph}。没有别名分析、没有布局传播、没有自动调整、没有 C++/CUDA/PTX 代码生成，也没有内存中 \texttt{Graph} 之外的持久格式。

\begin{table}[H]
\centering
\caption{Stax IR：存在的内容与故意缺席的内容。每个缺失的行都是 \texttt{Composite} 到 \texttt{tpx::ops} 的回退。}
\label{tab:staxscope}
\small
\begin{tabular}{p{3.6cm}p{4.4cm}p{3.8cm}}
\toprule
类别 & 目前在 \texttt{stax} & 故意缺席 \\
\midrule
红外节点 & \texttt{ValueNode} / \texttt{OpNode}，顺序 \texttt{execute} & 没有\texttt{Block}，没有\texttt{Function}，没有一流的控制流 \\
Ops & \texttt{div}, \texttt{t}, \texttt{relu}, \texttt{max\_pool2d}, \texttt{flatten} & 没有量化/稀疏/自定义调度键 \\
通行证 & 一次垂直逐点融合 (\texttt{fused\_mul\_add}) & 没有循环融合，没有布局规划器，没有矢量化策略 \\
后端 & 重新输入 \texttt{tpx::ops} 以便设备/区分保持正确 & 没有独立的代码生成器（\texttt{tensorplay.compile} 是条目） \\
定制操作 & \texttt{custom\_op} 节点 + \texttt{setCustomOpExecutor} 位于 \texttt{Graph.cpp} & 没有用于自定义操作的用户可见的 IR 构建器（使用 eager ops） \\
\bottomrule
\end{tabular}
\end{table}

目标消费者不是直接编写 \texttt{stax.Graph} 的 Python 作者，而是引用的更高级别条目 \texttt{tensorplay.compile(backend="stax")}。开发新融合模式原型的研究人员用 C++ 编写 \texttt{Graph}，通过 \texttt{REGISTER\_STAX\_PASS} 注册 \texttt{Pass}，并通过 \texttt{tensorplay.\_C.\_stax} 执行它，而无需更改面向用户的 \texttt{compile} 合约。


